\section{A variational principle}\label{sec:compactness}

A counterexample to $\Gtz(m,k)$ is a design whose selection objective
falls below the threshold.  We would like the worst one, and to read off
what minimality forces.  A worst design need not exist: the space of
designs at a given size and rank is not compact, and a minimizing
sequence need not converge in it (Example~\ref{ex:escape}).

Grant the conjecture at the two cells one size below.  Then a
counterexample at $(m,k)$ may be taken to minimize the chart objective
globally over a compact space, and to have strictly positive weights
(Theorem~\ref{thm:variational}).

The step that discharges the boundary holds at vanishing weight and
fails at every positive weight, however small (Theorem~\ref{thm:sharp}).
It fails on the extremal locus of \S\ref{sec:ties}.  So the argument runs
at the limit point.

\subsection{The compactification}\label{sec:compactchart}

Fix $m$ and $k$.  In a design the weights are confined to the open
simplex and the vectors are not confined.  Parseval controls only the
product: by Lemma~\ref{lem:trace},
\begin{equation}\label{eq:levbound}
  \sum_c s_c\;=\;\sum_c t_c\ell_c\;=\;k,
  \qquad\text{so}\qquad \ell_c\;\le\;\frac{k}{t_c} .
\end{equation}
As a weight tends to zero the corresponding vector may therefore grow
without bound, and it may do so while keeping its leverage score
$s_c=t_c\ell_c$ fixed.

\begin{example}\label{ex:escape}
Let $k=1$, $m=2$ and, for $0<u\le\tfrac12$,
\[
  t_1=1-u,\quad t_2=u,\qquad
  g_1=\frac1{\sqrt{2(1-u)}},\quad g_2=\frac1{\sqrt{2u}} .
\]
Then $t_1g_1^2+t_2g_2^2=\tfrac12+\tfrac12=1$, so this is a design at
$(2,1)$ for every such $u$.  Its leverages are
$\ell_1=\tfrac1{2(1-u)}$ and $\ell_2=\tfrac1{2u}$, and
$\ell_2\to\infty$ as $u\to0$.
No subsequence converges in the design space.  Its leverage scores are
$s_1=s_2=\tfrac12$, constant in $u$.
\end{example}

The chart of \S\ref{sec:chart} is the change of variables that
retains the leverage scores and discards the divergence.  With
$v_c=\sqrt{t_c}\,g_c$ and $P=VV\adj$ as in \eqref{eq:chart},
\begin{equation}\label{eq:pivotlev}
  P_{cc}\;=\;\lVert v_c\rVert^2\;=\;t_c\ell_c\;=\;s_c ,
\end{equation}
so the divergence in \eqref{eq:levbound} sits entirely in the factor
$1/t_c$: what the chart records is the leverage score, and that stays in
$[0,1]$ by Lemma~\ref{lem:share}.  Figure~\ref{fig:twosimplices} draws the
two coordinates at $(4,3)$.  The domination criterion, by
Lemma~\ref{lem:chart}, is a statement about $(P,t)$ alone.  In
Example~\ref{ex:escape} the chart is constant: $v_1=v_2=1/\sqrt2$, so
\[
  P=\begin{pmatrix}\tfrac12&\tfrac12\\[2pt]\tfrac12&\tfrac12\end{pmatrix},
  \qquad t=(1-u,\,u)\;\longrightarrow\;(1,0)\quad(u\to0),
\]
and the family converges, in these coordinates, to a pair $(P,t)$ whose
second weight vanishes while the corresponding diagonal entry of $P$ does
not.  Such pairs are the points to be added.

\begin{figure}[ht]
  \centering
  % The two simplices at (4,3): the weights, and the leverage scores through the
% substitution s_c -> 1 - s_c.  Both are 3-simplices inside R^4, drawn through
% the affine identification of the standard simplex with R^3 that carries the
% vertex e_c to the c-th vertex of the regular tetrahedron of
% figures/tetrahedron.tex,
%   x_1 = p_1 + p_2 - p_3 - p_4
%   x_2 = p_1 - p_2 + p_3 - p_4
%   x_3 = p_1 - p_2 - p_3 + p_4        (barycentre -> origin),
% followed by the orthographic projection of R^3 with azimuth 15 and elevation
% 70 degrees, scaled by 1.35cm -- the camera and the scale of
% figures/tetrahedron.tex, so the tetrahedron there and these two are the same
% drawing:
%   X = 1.35(-0.2588 x_1 + 0.9659 x_2)
%   Y = 1.35(-0.9077 x_1 - 0.2432 x_2 + 0.3420 x_3)
% The right panel repeats the left, shifted 7.7cm in X.  The leverage-score
% polytope is a simplex because the corank m - k is one; at corank two and above
% it is a hypersimplex, and this drawing does not apply.
%
% The path drawn is t = (1-3u, u, u, u) for 0 < u <= 1/4, carried at the chart
% P = I_4 - J_4/4 of the tetrahedron design: s_c = P_cc = 3/4 at every u, and
% l_c = s_c/t_c is 3/(4(1-3u)) at c = 1 and 3/(4u) at c = 2, 3, 4, all equal to
% 3 at u = 1/4.  In the identification above the path is the segment from the
% origin to the vertex e_1, since (1-3u)g_1 + u(g_2 + g_3 + g_4) = (1-4u)g_1.
\begin{tikzpicture}[
  x=1cm, y=1cm,
  vec/.style={-{Latex[length=2mm,width=1.5mm]}, line width=0.75pt},
  frame/.style={densely dotted, line width=0.3pt, black!35},
  lead/.style={line width=0.3pt, black!55},
  edge/.style={line width=0.45pt, black!65},
  hidden/.style={line width=0.45pt, black!45, dashed, dash pattern=on 1.4pt off 1.2pt},
  every node/.style={font=\small}]

% ---- one simplex, vertices e_1..e_4, barycentre at the origin -----------
% projected vertices:  e_1 -> ( 1, 1, 1) -> ( 0.9546,-1.0920)
%                      e_2 -> ( 1,-1,-1) -> (-1.6534,-1.3588)
%                      e_3 -> (-1, 1,-1) -> ( 1.6534, 0.4354)
%                      e_4 -> (-1,-1, 1) -> (-0.9546, 2.0154)
% e_2--e_3 is the far diagonal of the silhouette quadrilateral, hence hidden;
% the four dotted medians locate the barycentre against the vertices.
\newcommand{\simplexframe}{%
  \draw[hidden] (-1.6534,-1.3588)--( 1.6534, 0.4354);
  \draw[edge]   ( 0.9546,-1.0920)--(-1.6534,-1.3588)
                ( 0.9546,-1.0920)--( 1.6534, 0.4354)
                ( 0.9546,-1.0920)--(-0.9546, 2.0154)
                (-1.6534,-1.3588)--(-0.9546, 2.0154)
                ( 1.6534, 0.4354)--(-0.9546, 2.0154);
  \draw[frame]  (0,0)--( 0.9546,-1.0920) (0,0)--(-1.6534,-1.3588)
                (0,0)--( 1.6534, 0.4354) (0,0)--(-0.9546, 2.0154);
  \fill (0,0) circle (1.6pt);}

% ======================= left: the weights ==============================
\simplexframe
\node[anchor=south] at (0,2.24) {the weights};
\draw[lead] (-0.13,0)--(-1.38,0);
\node[anchor=east, align=right, font=\footnotesize] at (-1.42,0)
  {$t_c=\tfrac14$\\[1pt] {\scriptsize at $u=\tfrac14$}};
% the escape path t = (1-3u, u, u, u) is the median from the barycentre to e_1
\draw[vec] (0,0)--( 0.9546,-1.0920);
\node[anchor=north, font=\footnotesize] at (0.9546,-1.30) {$t=(1,0,0,0)$};
\node[anchor=north, align=center, font=\footnotesize] at (0,-1.90)
  {$t_c\ge0$, \ $\sum_c t_c=1$\\ $t\in\Delta^{3}$};

% ======================= middle: what each coordinate does ==============
\node[align=center, font=\footnotesize] at (3.85,0.34)
  {the chart $P$ held fixed\\[1pt]
   $t=(1-3u,u,u,u)$, \ $u\to0$\\[5pt]
   $s_c=P_{cc}$\\[5pt]
   $\ell_c=s_c/t_c$\\[1pt]
   {\scriptsize singular only at $t_c=0$}\\[5pt]
   $\ell_1\to\tfrac34$\\[1pt]
   $\ell_2=\ell_3=\ell_4=\tfrac{3}{4u}\to\infty$};

% ======================= right: the leverage scores =====================
\begin{scope}[shift={(7.7,0)}]
  \simplexframe
  \node[anchor=south] at (0,2.24) {the leverage scores};
  \draw[lead] (0.13,0)--(1.66,0);
  \node[anchor=west, align=left, font=\footnotesize] at (1.70,0)
    {$s_c=\tfrac34$\\[1pt] {\scriptsize at every $u$}};
  \node[anchor=north, font=\footnotesize] at (0.9546,-1.30) {$s=(0,1,1,1)$};
  \node[anchor=north, align=center, font=\footnotesize] at (0,-1.90)
    {$0\le s_c\le1$, \ $\sum_c s_c=3$\\ $1-s_c\in\Delta^{3}$};
\end{scope}

\end{tikzpicture}

  \caption{The two simplices at $(4,3)$.  Left: the weights, in
  $\Delta^{3}$.  Right: the leverage scores, confined by
  Lemmas~\ref{lem:trace} and~\ref{lem:share} to
  $\{s\in[0,1]^{4}:\sum_c s_c=3\}$, which at corank one is again
  $\Delta^{3}$, in the coordinates $1-s_c$.  The tetrahedron design of
  Example~\ref{ex:tetra} is the barycentre of each.  With its chart held
  fixed (Lemma~\ref{lem:realize}) the weights run to a vertex while the
  leverage scores stand still, and the leverages $\ell_c=s_c/t_c$
  diverge; Example~\ref{ex:escape} is the same motion at $(2,1)$.  The
  leverage score is the coordinate that stays compact and the weight the
  one that reaches the boundary, so the compactification is taken in
  $(P,t)$ rather than in the design data.}
  \label{fig:twosimplices}
\end{figure}

\begin{definition}\label{def:chartpoint}
A \emph{chart point} of size $m$ and rank $k$ is a pair $(P,t)$ in which
$P$ is a real symmetric idempotent $m\times m$ matrix with $\tr P=k$ and
$t$ lies in the closed simplex
$\Delta^{m-1}=\{t\in\R^m:t_c\ge0,\ \sum_ct_c=1\}$.  Put $D:=\diag(t)$,
$W:=P-D$, and
\begin{equation}\label{eq:objective}
  G(P,t)\;:=\;\max_{\lvert C\rvert=k}\;\lmin\bigl(W[C]\bigr) .
\end{equation}
A $k$-subset $C$ \emph{dominates} at $(P,t)$ if $W[C]\psd0$.  Write
$\Gtz^{\mathrm{ch}}(m,k)$ for the assertion that every chart point of
size $m$ and rank $k$ admits a dominating $k$-subset.
\end{definition}

Since a maximum over a finite set is attained,
\begin{equation}\label{eq:threshold}
  G(P,t)\ge0
  \quad\Longleftrightarrow\quad
  \text{some $k$-subset dominates at $(P,t)$.}
\end{equation}
No condition ties $P$ to $t$.

The relation \eqref{eq:pivotlev} is the definition of $s_c$, not a
constraint, and on the compactification the vector lengths are recovered
only where $t_c>0$, by $\ell_c=P_{cc}/t_c$.  That formula is the
boundary's only singularity.

\begin{example}\label{ex:tetragap}
The tetrahedron of Example~\ref{ex:tetra} has $P=I_4-\tfrac14J_4$ and
$t_c=\tfrac14$ by Example~\ref{ex:tetrachart}, so
\[
  W\;=\;\tfrac34I_4-\tfrac14J_4,
  \qquad
  W[C]\;=\;\tfrac14\bigl(3I_3-J_3\bigr)
\]
at each of the four triples.  Since $J_3$ has eigenvalues $3,0,0$, the
spectrum of $W[C]$ is $\{0,\tfrac34,\tfrac34\}$, and $G(P,t)=0$: every
triple sits on the threshold and none crosses it, which is
\eqref{eq:threshold} read at a design whose triples are known to dominate
without slack.  The gap $\SC-I_3$ has spectrum $\{0,3,3\}$ by
\eqref{eq:tetraspec}, four times this one.
\end{example}

\begin{proposition}\label{prop:chartdesign}
The chart points of size $m$ and rank $k$ with $t_c>0$ for every $c$ are
exactly the chart points of designs of size $m$ and rank $k$ over $\R$,
and the design is determined by the chart point up to an orthogonal
transformation of $\R^k$.  Under this correspondence a $k$-subset dominates in the design
if and only if it dominates at the chart point.
\end{proposition}

\begin{proof}
A design produces $(VV\tp,t)$, which satisfies
Definition~\ref{def:chartpoint}.  Here $VV\tp$ is symmetric; it is
idempotent because $VV\tp VV\tp=V(V\tp V)V\tp$ and $V\tp V=I_k$, the
columns of $V$ being orthonormal; and $\tr(VV\tp)=\tr(V\tp V)=k$ by
cyclicity of the trace.  The weights lie in the open simplex.
Conversely let $(P,t)$ be a chart point with all weights positive.  A
symmetric idempotent is the orthogonal projection onto its range, and
its eigenvalues are $0$ and $1$, so its rank and its trace agree:
$\rk P=\tr P=k$.  Let the columns of the $m\times k$ matrix $V$ be an
orthonormal basis of that range, so that $P=VV\tp$ and $V\tp V=I_k$.
Write $v_c\tp$ for the $c$-th row of $V$ and put $g_c:=v_c/\sqrt{t_c}$.
Then
\[
  \sum_c t_c\,g_c^{\vphantom{\mathsf{T}}}g_c\tp
  \;=\;\sum_c v_c^{\vphantom{\mathsf{T}}}v_c\tp\;=\;V\tp V\;=\;I_k ,
\]
which is \eqref{eq:parseval}; the weights are positive and sum to one; so
this is a design, and its chart point is $(P,t)$.  Two orthonormal bases
of the range differ by an orthogonal $Q$, and the substitution
$V\mapsto VQ$ sends $g_c\mapsto Q\tp g_c$ and $\SC\mapsto Q\tp\SC Q$.
The comparison \eqref{eq:dominate} with $I_k$ survives because
$Q\tp I_kQ=I_k$.  The last assertion is Lemma~\ref{lem:chart}.
\end{proof}

\begin{lemma}\label{lem:compact}
The set of chart points of size $m$ and rank $k$ is compact, and
nonempty when $k\le m$.
\end{lemma}

\begin{proof}
It is the product of $\Delta^{m-1}$, which is compact, with the set
$\Pi_k$ of symmetric idempotents of trace $k$.  The three conditions
defining $\Pi_k$ are polynomial equations in the entries, so $\Pi_k$ is
closed; and for $P\in\Pi_k$,
\[
  \lVert P\rVert_{F}^{2}\;=\;\tr(P\tp P)\;=\;\tr(P^{2})\;=\;\tr P\;=\;k ,
\]
so $\Pi_k$ lies on a sphere of the Frobenius norm and is bounded, hence
compact.  It is the Grassmannian of $k$-planes in $\R^m$, nonempty for
$k\le m$.
\end{proof}

\begin{lemma}\label{lem:continuity}
$G$ is continuous.  Indeed, for chart points $(P,t)$ and $(Q,t')$ of the
same size and rank,
\[
  \bigl\lvert G(P,t)-G(Q,t')\bigr\rvert
  \;\le\;\lVert P-Q\rVert_{2}+\max_c\lvert t_c-t'_c\rvert ,
\]
where $\lVert\cdot\rVert_{2}$ is the operator norm.
\end{lemma}

\begin{proof}
For symmetric $A$ one has
$\lmin(A)=\min_{\lVert x\rVert=1}x\tp Ax$, an infimum of functions of $A$
each of which satisfies
$\lvert x\tp Ax-x\tp Bx\rvert\le\lVert A-B\rVert_{2}$ for
$\lVert x\rVert=1$.  Two infima over the same index set differ by at
most the supremum of the difference, so
$\lvert\lmin(A)-\lmin(B)\rvert\le\lVert A-B\rVert_{2}$.  Restriction to
the principal block on an index set $C$ is $A\mapsto E\tp AE$, where the
columns of $E$ are the coordinate vectors $e_c$ with $c\in C$, so it
does not increase the operator norm.  The two gap matrices differ by
$(P-Q)-\diag(t-t')$, of operator norm at most
$\lVert P-Q\rVert_{2}+\max_c\lvert t_c-t'_c\rvert$.  So each of the
finitely many functions $(P,t)\mapsto\lmin(W[C])$ obeys the stated bound,
and so does their maximum.
\end{proof}

Chart points with positive weights are dense, by a mixing of the weights
that recurs twice below.

\begin{lemma}\label{lem:relax}
Let $(P,t)$ be a chart point of size $m$ and rank $k$, let
$\bar t:=(1/m,\dots,1/m)$ and, for $0\le\theta\le1$, put
$t^{\theta}:=(1-\theta)t+\theta\bar t$.  Then $(P,t^{\theta})$ is a chart
point, its weights are strictly positive for $\theta>0$, and for every
$x\in\R^m$
\begin{equation}\label{eq:relax}
  x\tp W^{\theta}x\;=\;x\tp Wx-\theta\sum_c
    \Bigl(\frac1m-t_c\Bigr)x_c^{2},
  \qquad W^{\theta}:=P-\diag(t^{\theta}) .
\end{equation}
Consequently the chart points with strictly positive weights are dense,
and
\[
  \min G\;=\;\inf\{\,G(P,t):\text{$(P,t)$ the chart point of a design}\,\}.
\]
\end{lemma}

\begin{proof}
The simplex is convex, so $t^{\theta}\in\Delta^{m-1}$, and
$t^{\theta}_c\ge\theta/m>0$ for $\theta>0$; $P$ is untouched.  Identity
\eqref{eq:relax} is
$\diag(t^{\theta})=\diag(t)+\theta\diag(\bar t-t)$ evaluated against $x$.
Letting $\theta\to0$ gives $t^{\theta}\to t$, which is the density
statement.  For the equality of extrema, $\min G$ exists by
Lemmas~\ref{lem:compact} and~\ref{lem:continuity}, and is at most the
infimum, a minimum over the chart points being at most an infimum over
some of them.  For the reverse, apply the mixing to an arbitrary chart
point $(A,u)$: at $\theta>0$ the weights of $u^{\theta}$ are strictly
positive, so $(A,u^{\theta})$ is the chart point of a design by
Proposition~\ref{prop:chartdesign}, and $G(A,u^{\theta})$ is at least
the infimum.  Lemma~\ref{lem:continuity} along $u^{\theta}\to u$ carries
that bound to $G(A,u)$, and in particular to the minimum.
\end{proof}

The design space itself admits no such statement.  It is the locus
$t_c>0$, which is open in the compactification and, by
Example~\ref{ex:escape}, not closed in it; a minimizing sequence of
designs need not converge to a design.

\subsection{The two objectives}\label{sec:objectives}

The objective \eqref{eq:objective} is not the design objective $F$ of
\eqref{eq:value}.  With
$D[C]=\diag(t_c)_{c\in C}$ the congruence \eqref{eq:congruence} reads
$W[C]=D[C]^{1/2}\bigl(\Gram_C-I_k\bigr)D[C]^{1/2}$ --- the Gram matrix,
not the subset sum --- so that, substituting $x=D[C]^{-1/2}y$,
\[
  \frac{x\tp W[C]x}{x\tp D[C]x}
  \;=\;\frac{y\tp(\Gram_C-I_k)y}{\lVert y\rVert^{2}} .
\]
For $\lvert C\rvert=k$ the matrix $M_C$ is square, so $\Gram_C=M_CM_C\tp$
and $\SC=M_C\tp M_C$ have the same characteristic polynomial and hence the
same least eigenvalue; minimizing the last display over $y\ne0$ therefore
gives $\lmin(\SC)-1$, and
\[
  F(\dsn)-1\;=\;\max_{\lvert C\rvert=k}\;\min_{x\ne0}
    \frac{x\tp W[C]x}{x\tp D[C]x},
\]
against $G=\max_C\min_{\lVert x\rVert=1}x\tp W[C]x$.

\begin{remark}\label{rem:objectives}
The two agree in sign, by \eqref{eq:threshold} and
Lemma~\ref{lem:chart}.  Beyond the sign they are related by a congruence
which is not a scalar unless the selected weights are equal, so they
need not share their minimizers.  The generalized form degenerates where
$D[C]$ becomes singular, that is on the points the compactification adds.
The absolute form does not.
\end{remark}

At a selection whose weights are all equal the congruence is that common
weight, and $\lmin(W[C])=t_c\bigl(\lmin(\SC)-1\bigr)$ at any $c\in C$.
Where the selected weights differ the factor differs with $C$, and the
order of two selections can invert.


Figure~\ref{fig:objectives} draws the two graphs and the six pairs.

\begin{figure}[ht]
  \centering
  % The design objective F against the chart objective G on one segment of
% weights at (4,2), and the six pairs read off at one point of that segment.
%
% PROJECTION.  There is none: the two drawn panels are graphs of functions of
% one real variable and the third panel is a table.  The two axis boxes are
% pgfplots boxes of exactly 5.8cm by 2.2cm (scale only axis), the upper with its
% south west corner at picture coordinate (0,0) and the lower at (0,-3.10cm), so
% the two share one abscissa and a vertical line at a given weight meets both at
% the same picture abscissa.  The maps from data to picture are
%   upper   X = 5.8 * t_0 / (2/3),   Y = 2.2 * (F-1) / 2.2
%   lower   X = 5.8 * t_0 / (2/3),   Y = 2.2 * G / 0.37
% and every plotted coordinate below is written in axis cs, that is as the pair
% (t_0, value), with its exact rational source in a trailing comment.  Decimals
% in the body are six-place roundings of those rationals and of nothing else:
%   4/15 = 0.266667   7/15 = 0.466667   8/15 = 0.533333   17/30 = 0.566667
%   3/5  = 0.6        19/30 = 0.633333  2/3  = 0.666667   1/6   = 0.166667
%   1/5  = 0.2        7/30 = 0.233333   1/3  = 0.333333
% The worst of them is off by 3.4e-7, which at 5.8cm over a range of 2/3
% displaces a point by 3e-6 cm.  The three marked weights are named outside both
% frames, at picture abscissae 5.8 * (8/15)/(2/3) = 4.64,
% 5.8 * (17/30)/(2/3) = 4.93 and 5.8 * (19/30)/(2/3) = 5.51: the two minimizers
% anchored south on the ordinate 2.30, clear of the upper frame, and the census
% weight anchored north on -0.23, centred in the 0.90-wide band that separates
% the two frames and that carries nothing else, the upper abscissa being left
% unlabelled.  The third panel sits at the abscissa 6.60, its two blocks anchored
% north west on the ordinates 2.42 and -1.60 with a rule between them on -1.42.
%
% THE CHART POINT, exactly.  On the index set {0,1,2,3},
%   P = [[4/5,2/5,0,0],[2/5,1/5,0,0],[0,0,1/2,1/2],[0,0,1/2,1/2]],
%   t = ( t_0, 2/3 - t_0, 1/6, 1/6 ),   0 <= t_0 <= 2/3 .
% The two diagonal blocks are u u^T/5 with u = (2,1) and w w^T/2 with w = (1,1);
% since ||u||^2 = 5 and ||w||^2 = 2 each is a rank-one orthogonal projection, so
% P^T = P, P^2 = P and tr P = 1 + 1 = 2.  The weights are nonnegative and sum to
% t_0 + (2/3 - t_0) + 1/6 + 1/6 = 1.  So (P,t) is a chart point of size four and
% rank two at every t_0, and the chart point of a design at every interior t_0.
%
% THE DESIGN, exactly.  The range of P is spanned by (2,1,0,0)/sqrt5 and
% (0,0,1,1)/sqrt2, so the rows of V are
%   v_0 = (2/sqrt5,0)  v_1 = (1/sqrt5,0)  v_2 = (0,1/sqrt2)  v_3 = (0,1/sqrt2)
% and g_c = v_c/sqrt(t_c) puts g_0, g_1 on the first axis of R^2 and g_2, g_3 on
% the second, with leverages and leverage scores
%   ell_0 = 4/(5 t_0)          s_0 = P_00 = 4/5
%   ell_1 = 1/(5(2/3 - t_0))   s_1 = P_11 = 1/5
%   ell_2 = ell_3 = 3          s_2 = s_3 = 1/2 .
% Parseval holds over the rationals at every interior t_0: sum_c t_c g_c g_c^T is
% diag(s_0 + s_1, s_2 + s_3) = diag(1,1) = I_2, and sum_c s_c = 2 = k.
%
% THE TWO OBJECTIVES.  For a mixed pair P_ab = 0, so W[C] = diag(A_a,A_b) with
% A_c = s_c - t_c, and Gram_C = diag(ell_a, ell_b).  A pair drawn from one block
% has parallel vectors, so there its Gram matrix is singular; sharing a
% characteristic polynomial with S_C, and S_C being positive semidefinite, that
% puts lmin(S_C) = 0.  And det(Gram_C - I_2) = -(ell_a + ell_b - 1) < 0, because
% ell_0 >= 6/5 at every interior t_0 and ell_2 = ell_3 = 3, so det W[C] < 0 by
% cor:minor.  (ell_1 is NOT bounded below by 6/5.  It rises from 3/10 at t_0 = 0
% and passes 6/5 only at t_0 = 1/2, so the bound must be carried by ell_0 alone.)
% The mixed pair {0,2} carries lmin(S_C) = min(ell_0,3) >= 6/5 and
% lmin(W[C]) = min(A_0,1/3) >= 2/15, above both, so neither maximum is attained
% inside a block.  The two indices of a mixed pair range independently over the
% blocks, so a maximum over mixed pairs of a minimum is the minimum of the two
% block maxima:
%   F = min( max(ell_0, ell_1), 3 ),      G = min( max(A_0, A_1), 1/3 ) .
% Here A_0 = 4/5 - t_0 and A_1 = t_0 - 7/15, of constant sum 1/3, so max(A_0,A_1)
% is a V with vertex at t_0 = 19/30 and value 1/6; and ell_0 falls while ell_1
% rises, so max(ell_0,ell_1) has its least value 3/2 at t_0 = 8/15.  Hence
% min G = 1/6 at t_0 = 19/30 and min F = 3/2 at t_0 = 8/15, each at that one
% weight and nowhere else.  Breakpoints of the two graphs:
%   F - 1 = 2 on (0,4/15];  4/(5t_0) - 1 on [4/15,8/15];
%           1/(5(2/3-t_0)) - 1 on [8/15,3/5];  2 on [3/5,2/3)
%   G     = 1/3 on [0,7/15];  4/5 - t_0 on [7/15,19/30];  t_0 - 7/15 on [19/30,2/3]
% with F-1 = 1/2 at 8/15, = 1 at 17/30, = 2 at 19/30, and G = 4/15 at 8/15,
% = 7/30 at 17/30, = 1/6 at 19/30, = 1/3 at 0, = 1/5 at 2/3.  On the whole
% segment F >= 3/2 and G >= 1/6, so the family carries no counterexample, and
% none is claimed: (4,2) is corank two, where the conjecture is a theorem.
%
% F IS A DESIGN OBJECTIVE, so it has no value at a weight no design charts, and
% the two endpoints of the segment are such weights.  Both horizontal segments
% are drawn closed all the same: at either end max(ell_0,ell_1) diverges while
% the truncating constant ell_2 = 3 binds, so F extends continuously with value
% 3.  G is defined at both endpoints, being a function of the chart point alone.
%
% THE CENSUS ABSCISSA.  t_0 = 17/30 lies in the open interval
% (8/15, 19/30) between the two minimizers, which is where the leverage maximum
% of the first block has moved to the index 1 while its slack maximum still sits
% at the index 0; the two objectives therefore select different pairs there, and
% 17/30 keeps every entry rational or a single surd.  At it t = (17,3,5,5)/30 and
% ell = (24/17, 2, 3, 3), and, least eigenvalue first,
%          spec(Gram_C - I_2)   spec W[C]
%   {0,1}  -1, 41/17            1/6 - sqrt37/15, 1/6 + sqrt37/15
%   {0,2}  7/17, 2              7/30, 1/3
%   {0,3}  7/17, 2              7/30, 1/3
%   {1,2}  1, 2                 1/10, 1/3
%   {1,3}  1, 2                 1/10, 1/3
%   {2,3}  -1, 5                -1/6, 5/6
% The {0,1} row: the trace is A_0 + A_1 = 1/3 and the determinant is
% (7/30)(1/10) - (2/5)^2 = -41/300, so the eigenvalues are
% (1/3 +- sqrt(1/9 + 41/75))/2 with 1/9 + 41/75 = 148/225.  The {2,3} row is the
% one whose selected weights agree: D[C] = I_2/6 there, and spec W[C] is
% spec(Gram_C - I_2) divided by 6, entry for entry.
%
% DISCLOSED.  NOTHING IN THIS FIGURE IS MECHANIZED.  Neither F nor G is a Lean
% definition -- the compactness development carries the predicate
% Gtz.ChartDominates and no real-valued chart objective -- and no declaration
% evaluates either at this chart point or exhibits this design.  The results the
% panels invoke (lem:congruence, lem:chart, cor:minor, thm:boundary) carry their
% own tags; the arithmetic here carries none.  Nearest miss, recorded so that a
% later reader does not take the absence for an oversight: the anonymous function
% sup'_{C in gates} lambdaMinMat(subsetSumRaw config C) - 1 appears inside
% Gtz.continuous_dominationMargin, and that is F - 1 in all but name.  It is not
% a definition, its maximum runs over a supplied gate family rather than over all
% k-subsets, its argument is an unconstrained pair (vectors, weights) with no
% Parseval condition, and nothing evaluates it anywhere.
\begin{tikzpicture}[
  gcurve/.style={line width=1.1pt},
  guide/.style={line width=0.5pt, black!55, dashed,
                dash pattern=on 2.2pt off 1.6pt},
  census/.style={line width=0.4pt, black!45, densely dotted},
  hair/.style={line width=0.4pt, black!45}]

% ================= upper panel: the design objective F - 1 =================
\begin{axis}[
  name=Fax, at={(0,0)},
  width=5.8cm, height=2.2cm, scale only axis,
  xmin=0, xmax=0.666667, ymin=0, ymax=2.2,
  xtick={0,0.2,0.4,0.6}, xticklabel=\empty,
  ytick={0,1,2}, yticklabels={$0$,$1$,$2$},
  yticklabel style={font=\scriptsize},
  ylabel={$F(\dsn)-1$}, ylabel style={font=\footnotesize},
  tick style={line width=0.4pt}, axis line style={line width=0.5pt},
  clip=true]

\draw[guide]  (axis cs:0.533333,0) -- (axis cs:0.533333,2.2);   % t_0 = 8/15
\draw[guide]  (axis cs:0.633333,0) -- (axis cs:0.633333,2.2);   % t_0 = 19/30
\draw[census] (axis cs:0.566667,0) -- (axis cs:0.566667,2.2);   % t_0 = 17/30

\draw[gcurve] (axis cs:0,2)   -- (axis cs:0.266667,2);          % (0,2)--(4/15,2)
\draw[gcurve] (axis cs:0.6,2) -- (axis cs:0.666667,2);          % (3/5,2)--(2/3,2)
\addplot[gcurve, domain=0.266667:0.533333, samples=90]          % 4/(5t_0) - 1
  {4/(5*x) - 1};
\addplot[gcurve, domain=0.533333:0.6, samples=60]               % 1/(5(2/3-t_0)) - 1
  {1/(5*(0.666667 - x)) - 1};

\fill (axis cs:0.533333,0.5) circle (1.7pt);                    % (8/15, 1/2)
\draw[line width=0.4pt, fill=white]
      (axis cs:0.566667,1) circle (1.5pt)                       % (17/30, 1)
      (axis cs:0.633333,2) circle (1.5pt);                      % (19/30, 2)
\node[font=\scriptsize, anchor=west, inner sep=2pt]
  at (axis cs:0.024,1.44) {$\min F=\tfrac32$};
\end{axis}

% ================= lower panel: the chart objective G ======================
\begin{axis}[
  name=Gax, at={(0,-3.10cm)},
  width=5.8cm, height=2.2cm, scale only axis,
  xmin=0, xmax=0.666667, ymin=0, ymax=0.37,
  xtick={0,0.2,0.4,0.6}, xticklabels={$0$,$\tfrac15$,$\tfrac25$,$\tfrac35$},
  xticklabel style={font=\scriptsize},
  ytick={0,0.166667,0.333333}, yticklabels={$0$,$\tfrac16$,$\tfrac13$},
  yticklabel style={font=\scriptsize},
  xlabel={$t_0$}, ylabel={$G(P,t)$},
  xlabel style={font=\footnotesize}, ylabel style={font=\footnotesize},
  tick style={line width=0.4pt}, axis line style={line width=0.5pt},
  clip=true]

\draw[guide]  (axis cs:0.533333,0) -- (axis cs:0.533333,0.37);  % t_0 = 8/15
\draw[guide]  (axis cs:0.633333,0) -- (axis cs:0.633333,0.37);  % t_0 = 19/30
\draw[census] (axis cs:0.566667,0) -- (axis cs:0.566667,0.37);  % t_0 = 17/30

\draw[gcurve] (axis cs:0,0.333333)                              % (0, 1/3)
           -- (axis cs:0.466667,0.333333)                       % (7/15, 1/3)
           -- (axis cs:0.633333,0.166667)                       % (19/30, 1/6)
           -- (axis cs:0.666667,0.2);                           % (2/3, 1/5)

\fill (axis cs:0.633333,0.166667) circle (1.7pt);               % (19/30, 1/6)
\draw[line width=0.4pt, fill=white]
      (axis cs:0.533333,0.266667) circle (1.5pt)                % (8/15, 4/15)
      (axis cs:0.566667,0.233333) circle (1.5pt);               % (17/30, 7/30)
\node[font=\scriptsize, anchor=west, inner sep=2pt]
  at (axis cs:0.024,0.070) {$\min G=\tfrac16$};
\end{axis}

% ---- the three marked weights, named outside the frames ------------------
\node[font=\scriptsize, anchor=south] at (4.64,2.30) {$\tfrac8{15}$};
\node[font=\scriptsize, anchor=south] at (5.51,2.30) {$\tfrac{19}{30}$};
\node[font=\scriptsize, anchor=north, inner sep=1pt]
  at (4.93,-0.23) {$\tfrac{17}{30}$};

% ============ the six pairs at t_0 = 17/30, and the boundary ==============
\node[anchor=north west, align=left, font=\scriptsize, inner sep=0pt]
  at (6.60,2.42) {%
  $m=4$, \ $k=2$, \ $t=\tfrac1{30}(17,3,5,5)$, \
  $\ell=(\tfrac{24}{17},2,3,3)$\\[4pt]
  \begin{tabular}{@{}l@{\hspace{1.4em}}l@{\hspace{1.4em}}l@{\hspace{0.8em}}l@{}}
    $C$ & $\operatorname{spec}(\Gram_C-I_2)$ & $\operatorname{spec}W[C]$ &\\[1.5pt]
    $\{0,1\}$ & $-1,\ \tfrac{41}{17}$ & $\tfrac16\mp\tfrac1{15}\sqrt{37}$ &\\[1.5pt]
    $\{0,2\}$ & $\tfrac7{17},\ 2$ & $\tfrac7{30},\ \tfrac13$ & $\leftarrow G$\\[1.5pt]
    $\{0,3\}$ & $\tfrac7{17},\ 2$ & $\tfrac7{30},\ \tfrac13$ & $\leftarrow G$\\[1.5pt]
    $\{1,2\}$ & $1,\ 2$ & $\tfrac1{10},\ \tfrac13$ & $\leftarrow F$\\[1.5pt]
    $\{1,3\}$ & $1,\ 2$ & $\tfrac1{10},\ \tfrac13$ & $\leftarrow F$\\[1.5pt]
    $\{2,3\}$ & $-1,\ 5$ & $-\tfrac16,\ \tfrac56$ &\\
  \end{tabular}};

\draw[hair] (6.60,-1.42) -- (13.55,-1.42);

\node[anchor=north west, align=left, font=\scriptsize, inner sep=0pt]
  at (6.60,-1.60) {%
  \begin{tabular}{@{}l@{\hspace{1.2em}}l@{\hspace{1.2em}}l@{}}
    at $\{0,2\}$ & $D[C]=\diag(\tfrac{17}{30},\tfrac16)$
      & $\tfrac{17}{30}\cdot\tfrac7{17}=\tfrac7{30}$\\[2pt]
    at $\{1,2\}$ & $D[C]=\diag(\tfrac1{10},\tfrac16)$
      & $\tfrac1{10}\cdot1=\tfrac3{30}$\\[2pt]
    at $\{2,3\}$ & $D[C]=\tfrac16I_2$
      & $\tfrac16\cdot(-1)=-\tfrac16$\\
  \end{tabular}\\[5pt]
  at $t_0=0$: \ $D[\{0,2\}]=\diag(0,\tfrac16)$ is singular,\\
  \phantom{at $t_0=0$: \ }%
  $e_0\tp W[\{0,2\}]e_0=\tfrac45$ and $\ell_0=4/(5t_0)$ diverges,\\
  \phantom{at $t_0=0$: \ }%
  while $W[\{0,2\}]=\diag(\tfrac45,\tfrac13)$ and $G=\tfrac13$};

\end{tikzpicture}

  \caption{The two objectives at $(4,2)$, on the weight segment of
  Example~\ref{ex:objectives}.  Left: the design objective above the
  chart objective, on one abscissa; the marked weight $\tfrac8{15}$
  carries the least value of the first and $\tfrac{19}{30}$ that of the
  second, and at $\tfrac{19}{30}$ the first is $3$ while at $\tfrac8{15}$
  the second is $\tfrac4{15}$.  Right: the six pairs at the dotted weight
  $\tfrac{17}{30}$, with an arrow on the pair each objective selects.
  The congruence \eqref{eq:congruence} makes the six rows agree in sign,
  by Sylvester's law, and lets them disagree in order: it carries the
  relative slack $\tfrac7{17}$ of $\{0,2\}$ up to $\tfrac7{30}$ and the
  larger relative slack $1$ of $\{1,2\}$ down to $\tfrac3{30}$, being
  a scalar only at $\{2,3\}$, whose two weights agree.  At the endpoint
  $t_0=0$ the leverage $\ell_0$ diverges and the generalized quotient on
  $\{0,2\}$ is unbounded above, while the gap stays finite and $\{0,2\}$
  dominates.  Neither
  objective is a Lean definition, and none of the arithmetic here is
  mechanized.}
  \label{fig:objectives}
\end{figure}

\subsection{The boundary}\label{sec:boundary}

Let $(P,t)$ be a chart point with $t_z=0$ for some index $z$.  The
diagonal entry $P_{zz}$ decides which of two identities disposes of it,
and that entry is a squared length.

\begin{lemma}[Lean]\label{lem:pivot}
For every chart point and every index $z$,
\begin{equation}\label{eq:pivot}
  \lVert Pe_z\rVert^{2}\;=\;e_z\tp P\tp Pe_z\;=\;e_z\tp P^{2}e_z
  \;=\;e_z\tp Pe_z\;=\;P_{zz} .
\end{equation}
Hence $P_{zz}\ge0$; and $P_{zz}=0$ forces $Pe_z=0$, that is, the $z$-th
column and, by symmetry, the $z$-th row of $P$ vanish identically.
\end{lemma}

\begin{proof}
The chain \eqref{eq:pivot} uses symmetry and idempotency in that order.
A vanishing sum of squares $\sum_cP_{cz}^{2}=P_{zz}=0$ has all terms
zero.
\end{proof}

By \eqref{eq:pivotlev} the quantity $P_{zz}$ is the limit of the leverage
score $s_z$.

The dichotomy is therefore the distinction between the two ways an atom
can leave a design: its Parseval mass vanishes, or its mass survives
while its weight vanishes and its vector escapes to infinity.
Example~\ref{ex:escape} is of the second kind, with $P_{zz}=\tfrac12$.

\subsubsection*{Case A: the atom vanishes}

\begin{lemma}[Lean]\label{lem:delete}
Let $(P,t)$ be a chart point of size $m$ and rank $k$ and let $z$ satisfy
$t_z<1$ and $Pe_z=0$.  Delete the index $z$: let $\widehat P$ be the
principal submatrix of $P$ on the remaining indices and
$\widehat t_c:=t_c/(1-t_z)$.  Then $(\widehat P,\widehat t)$ is a chart
point of size $m-1$ and rank $k$, and for every $y\in\R^m$ with $y_z=0$,
writing $\widehat y$ for its restriction,
\begin{equation}\label{eq:reweight}
  y\tp Wy\;=\;\widehat y\tp\widehat W\widehat y
    +\frac{t_z}{1-t_z}\sum_{c\ne z}t_cy_c^{2}
  \;\ge\;\widehat y\tp\widehat W\widehat y .
\end{equation}
Consequently a subset dominating at $(\widehat P,\widehat t)$ dominates,
under the inclusion of index sets, at $(P,t)$.
\end{lemma}

\begin{proof}
Symmetry is inherited.  Since the $z$-th row and column of $P$ vanish,
the deleted index contributes nothing to any product, so
$\widehat P\widehat P$ is the corresponding submatrix of $P^{2}=P$, and
$\tr\widehat P=\tr P-P_{zz}=k$.  The weights $\widehat t_c$ are
nonnegative and sum to $(1-t_z)/(1-t_z)=1$.  For the identity, $y_z=0$
gives $y\tp Py=\widehat y\tp\widehat P\widehat y$, while
\[
  \widehat y\tp\diag(\widehat t)\widehat y-y\tp Dy
  \;=\;\sum_{c\ne z}\Bigl(\frac{t_c}{1-t_z}-t_c\Bigr)y_c^{2}
  \;=\;\frac{t_z}{1-t_z}\sum_{c\ne z}t_cy_c^{2},
\]
which is \eqref{eq:reweight} after subtraction.  The right-hand side is
nonnegative because $t_z\ge0$ and $t_z<1$.
\end{proof}

The reweighting is favourable at every admissible $t_z$, not only at
$t_z=0$; the hypothesis $t_z=0$ is needed only at the pivot of the next
case, where \S\ref{sec:sharp} locates the failure.

\subsubsection*{Case B: the direction escapes}

\begin{lemma}[Lean]\label{lem:deflate}
Let $(P,t)$ be a chart point of size $m$ and rank $k\ge1$ and let
$P_{zz}>0$.  Then
\begin{equation}\label{eq:deflate}
  P'\;:=\;P-\frac{(Pe_z)(Pe_z)\tp}{P_{zz}}
\end{equation}
is symmetric and idempotent with $\tr P'=k-1$ and $P'e_z=0$.  Hence
$(P',t)$ is a chart point of size $m$ and rank $k-1$ whose $z$-th row and
column vanish.
\end{lemma}

\begin{proof}
Write $r:=Pe_z$ and $R:=rr\tp$, so that $P'=P-R/P_{zz}$; symmetry is
clear.  Idempotency of $P$ gives $Pr=P^{2}e_z=Pe_z=r$, hence $PR=R$ and,
transposing, $RP=R$.  Lemma~\ref{lem:pivot} gives
$RR=r(r\tp r)r\tp=P_{zz}R$.  Therefore
\[
  \Bigl(P-\frac{R}{P_{zz}}\Bigr)^{2}
  \;=\;P^{2}-\frac{PR+RP}{P_{zz}}+\frac{RR}{P_{zz}^{2}}
  \;=\;P-\frac{2R}{P_{zz}}+\frac{R}{P_{zz}}
  \;=\;P-\frac{R}{P_{zz}} .
\]
The trace drops by $\tr R/P_{zz}=\lVert r\rVert^{2}/P_{zz}=1$, again by
\eqref{eq:pivot}.  Finally $r\tp e_z=P_{zz}$, so
$P'e_z=r-r\,P_{zz}/P_{zz}=0$.
\end{proof}

Deflation compresses the surviving vectors onto the orthogonal complement
of the escaping direction, carried out in the chart: on the surviving
indices $P'$ is the Schur complement of the pivot $P_{zz}$ in $P$.  The
computation chooses no basis of that complement and takes no square root.
It performs one division, by the pivot.  Since $P=P'+rr\tp/P_{zz}$ with
$r=Pe_z$, for every $x\in\R^m$
\begin{equation}\label{eq:transverse}
  x\tp Wx\;=\;x\tp(P'-D)x+\frac{\langle Pe_z,x\rangle^{2}}{P_{zz}} ,
\end{equation}
so the gap form upstairs exceeds the deflated one by the rank-one pivot
term.  Splitting the probe along the escaping coordinate isolates that
term.

\begin{lemma}\label{lem:schur}
Let $(P,t)$ be a chart point with $P_{zz}>0$ and let $P'$ be as in
\eqref{eq:deflate}.  Write $x=y+x_ze_z$ with $y_z=0$.  Then
\begin{equation}\label{eq:schur}
  x\tp Wx\;=\;y\tp(P'-D)\,y
    \;+\;\frac{\bigl(\langle Pe_z,y\rangle+x_zP_{zz}\bigr)^{2}}{P_{zz}}
    \;-\;t_z\,x_z^{2} .
\end{equation}
\end{lemma}

\begin{proof}
Expand $x\tp Px$ against the coordinate spike.  With
$a:=\langle Pe_z,y\rangle$ and $e_z\tp Pe_z=P_{zz}$,
\[
  x\tp Px\;=\;y\tp Py+2x_za+x_z^{2}P_{zz},
\]
and $P=P'+rr\tp/P_{zz}$ gives
$y\tp Py=y\tp P'y+a^{2}/P_{zz}$.  Since $y_z=0$,
$x\tp Dx=y\tp Dy+t_zx_z^{2}$.  Subtracting,
\[
  x\tp Wx\;=\;y\tp(P'-D)y
    +\frac{a^{2}}{P_{zz}}+2x_za+x_z^{2}P_{zz}-t_zx_z^{2} .
\]
The three middle terms complete a square: since $P_{zz}>0$,
\[
  \frac{a^{2}}{P_{zz}}+2x_za+x_z^{2}P_{zz}
  \;=\;\frac{a^{2}+2x_zaP_{zz}+x_z^{2}P_{zz}^{2}}{P_{zz}}
  \;=\;\frac{\bigl(a+x_zP_{zz}\bigr)^{2}}{P_{zz}} ,
\]
which is the middle term of \eqref{eq:schur}.
\end{proof}

Identity \eqref{eq:schur} is exact at every weight.  Its three terms are,
in order: the deflated gap, a square, and a deficit proportional to the
weight at the pivot.  At $t_z=0$ the deficit is absent and the lift
follows.

\begin{theorem}[Lean]\label{thm:schurstep}
Let $(P,t)$ be a chart point of size $m$ and rank $k\ge1$, let $z$
satisfy $P_{zz}>0$ and $t_z=0$, and let $C$ be a subset of the remaining
indices dominating at the chart point $(P',t)$ of rank $k-1$.  Then
$C\cup\{z\}$ dominates at $(P,t)$.
\end{theorem}

\begin{proof}
Let $x$ be supported on $C\cup\{z\}$ and split $x=y+x_ze_z$ with $y_z=0$;
then $y$ is supported on $C$, so $y\tp(P'-D)y\ge0$ by hypothesis.  In
\eqref{eq:schur} the last term vanishes because $t_z=0$, and the middle
term is nonnegative because $P_{zz}>0$.  Hence $x\tp Wx\ge0$.
\end{proof}

The pivot matches because the weight there is zero: the deflation
\eqref{eq:deflate} subtracts $P_{zz}$ from the diagonal at $z$, and
\eqref{eq:schur} restores the same amount.  The lift consumes the
rank-one excess of \eqref{eq:transverse} in full, so none of it is slack.

\begin{example}\label{ex:escapeb}
Return to Example~\ref{ex:escape}.  At its limit point, $z=2$,
$P_{zz}=\tfrac12$
and $t_z=0$.  Here $Pe_2=(\tfrac12,\tfrac12)\tp$ with
$\lVert Pe_2\rVert^{2}=\tfrac12=P_{22}$, as \eqref{eq:pivot} requires,
and \eqref{eq:deflate} gives $P'=P-2\cdot\tfrac14
\left(\begin{smallmatrix}1&1\\1&1\end{smallmatrix}\right)=0$, of trace
$0=k-1$.  Deleting $z$ leaves the chart point of size $1$ and rank $0$,
at which the empty selection dominates; Theorem~\ref{thm:schurstep}
returns $\{2\}$, and indeed $W_{22}=P_{22}-t_2=\tfrac12\ge0$.
\end{example}

\subsubsection*{The dichotomy}

\begin{theorem}[Lean]\label{thm:boundary}
Let $m\ge2$ and $k\ge1$, and assume $\Gtz^{\mathrm{ch}}(m-1,k)$ and
$\Gtz^{\mathrm{ch}}(m-1,k-1)$.  Then every chart point of size $m$ and
rank $k$ with $t_z=0$ for some $z$ admits a dominating $k$-subset.
\end{theorem}

\begin{proof}
By Lemma~\ref{lem:pivot} either $P_{zz}=0$ or $P_{zz}>0$.

If $P_{zz}=0$ then $Pe_z=0$, and $t_z=0<1$, so Lemma~\ref{lem:delete}
produces a chart point of size $m-1$ and rank $k$.  By
$\Gtz^{\mathrm{ch}}(m-1,k)$ some $k$-subset dominates there, and by
\eqref{eq:reweight} its image dominates at $(P,t)$.

If $P_{zz}>0$, deflate: by Lemma~\ref{lem:deflate}, $(P',t)$ is a chart
point of rank $k-1$ whose $z$-th column vanishes, so
Lemma~\ref{lem:delete} applies to it and produces a chart point of size
$m-1$ and rank $k-1$.  By $\Gtz^{\mathrm{ch}}(m-1,k-1)$ some
$(k-1)$-subset $C$ dominates there; by \eqref{eq:reweight} it dominates
at $(P',t)$; and by Theorem~\ref{thm:schurstep}, $C\cup\{z\}$, of
cardinality $k$, dominates at $(P,t)$.
\end{proof}

Both branches reduce the size by one, so an induction on the dichotomy is
well founded on the size alone.

At $k=1$ the second hypothesis is
$\Gtz^{\mathrm{ch}}(m-1,0)$, which holds because the empty selection
dominates.  Every step of the reduction is an identity, so it consumes no
slack and stays compatible with Corollary~\ref{cor:nomargin}.
Figure~\ref{fig:boundary} separates the two branches.

\begin{corollary}[Lean]\label{cor:interiorcex}
Under the hypotheses of Theorem~\ref{thm:boundary}, every chart point of
size $m$ and rank $k$ admitting \emph{no} dominating $k$-subset has
$t_c>0$ for every $c$.
\end{corollary}

\begin{proof}
Contrapositive of Theorem~\ref{thm:boundary}.
\end{proof}

\begin{figure}[ht]
  \centering
  \begin{tikzpicture}[
  font=\small,
  bx/.style={draw, rounded corners=1pt, inner sep=3pt, minimum height=5mm,
             minimum width=17mm, align=center, font=\small},
  ar/.style={-{Latex[length=1.7mm]}, line width=0.5pt},
  ttl/.style={font=\scriptsize\itshape, align=center},
  txt/.style={font=\scriptsize, align=center, inner sep=1.5pt}]

% =====================================================================
% (a)  the compactification:  P fixed, t in the closed simplex
% =====================================================================
\begin{scope}
\coordinate (v1) at (0,0);
\coordinate (v2) at (4.6,0);
\coordinate (v3) at (2.3,3.98);
\coordinate (w1) at (0.42,0.243);
\coordinate (w2) at (4.18,0.243);
\coordinate (w3) at (2.3,3.494);

\fill[black!7] (w1) -- (w2) -- (w3) -- cycle;
\fill[pattern=north east lines, pattern color=black!45, even odd rule]
  (v1) -- (v2) -- (v3) -- cycle   (w1) -- (w2) -- (w3) -- cycle;
\draw[line width=1pt] (v1) -- (v2) -- (v3) -- cycle;
\draw[black!35, line width=0.3pt] (w1) -- (w2) -- (w3) -- cycle;
\foreach \p in {(v1),(v2),(v3)}{\fill \p circle (1.5pt);}

\node[txt] at (2.3,2.25) {interior\\ $t_c>0$};
\fill (2.3,1.42) circle (1.7pt);
\node[txt, anchor=north] at (2.3,1.26) {$(P_{\min},t_{\min})$};
\node[txt, anchor=north] at (2.3,-0.06) {$t_z=0$};

\node[txt, anchor=north] at (2.3,-0.52)
  {closed simplex $\Delta^{m-1}$: $t_c\ge0$, $\sum_c t_c=1$;\\[1pt]
   the hatched faces $t_c=0$ carry no counterexample,\\[1pt]
   so a minimiser of $G$ lies in the interior};
\node[txt, anchor=north] at (2.3,-1.72)
  {$G(P,t)=\displaystyle\max_{|C|=k}\lmin\bigl(W[C]\bigr)$,\quad $W=P-\diag t$,\\[2pt]
   continuous on a compact set; the minimum is attained};
\end{scope}

% =====================================================================
% (b)  the two boundary mechanisms at a vanishing weight
% =====================================================================
\begin{scope}[xshift=5.9cm]
\node[txt, anchor=north west, text width=79mm, align=left] at (0,5.05)
  {A chart point at $(m,k)$ with $t_z=0$. The pivot identity
   $\lVert Pe_z\rVert^{2}=P_{zz}$ splits two cases.};

% ---- case A : the atom vanishes -------------------------------------
\node[ttl] at (1.95,4.05) {atom deletion};
\node[bx] (a1) at (1.95,3.62) {$P_{zz}=0$};
\node[txt, text width=38mm, align=center] (a2) at (1.95,2.52)
  {the column $Pe_z$ vanishes;\\ delete $z$, rescale $t\mapsto t/(1-t_z)$};
\node[bx] (a3) at (1.95,1.40) {$(m-1,k)$};
\node[txt, text width=38mm, align=center] (a4) at (1.95,0.78) {$C$ lifts verbatim};
\node[bx] (a5) at (1.95,0.15) {$|C|=k$};
\draw[ar] (a1) -- (a2); \draw[ar] (a2) -- (a3);
\draw[ar] (a3) -- (a4); \draw[ar] (a4) -- (a5);

% ---- case B : deflation at the pivot --------------------------------
\node[ttl] at (6.05,4.05) {Schur deflation};
\node[bx] (b1) at (6.05,3.62) {$P_{zz}>0$};
\node[txt, text width=38mm, align=center] (b2) at (6.05,2.52)
  {$P'=P-(Pe_z)(Pe_z)\tp/P_{zz}$,\\ $\tr P'=k-1$, $P'e_z=0$; delete $z$};
\node[bx] (b3) at (6.05,1.40) {$(m-1,k-1)$};
\node[txt, text width=38mm, align=center] (b4) at (6.05,0.78) {$C\mapsto C\cup\{z\}$};
\node[bx] (b5) at (6.05,0.15) {$|C\cup\{z\}|=k$};
\draw[ar] (b1) -- (b2); \draw[ar] (b2) -- (b3);
\draw[ar] (b3) -- (b4); \draw[ar] (b4) -- (b5);

% ---- the identity that pays for the second case ---------------------
\node[txt, anchor=north] at (4.0,-0.45)
  {$x\tp Wx=x\tp(P'-\diag t)\,x
    +\dfrac{\langle Pe_z,x\rangle^{2}}{P_{zz}}$,\\[3pt]
   and at $t_z=0$ the three terms complete a square};
\end{scope}

\end{tikzpicture}

  \caption{The compactification and the boundary dichotomy.  Left: the
  objective $G$ of \eqref{eq:objective} is continuous on a compact space,
  so it attains its minimum; Theorem~\ref{thm:boundary} empties the
  faces $t_z=0$ of counterexamples, which places a negative minimum in
  the interior.  Right: the two mechanisms, separated by the pivot
  identity \eqref{eq:pivot}.  At $P_{zz}=0$ the atom is deleted and the
  weights rescaled; at $P_{zz}>0$ the point is deflated by
  \eqref{eq:deflate} and the selected subset regains $z$, the deficit in
  \eqref{eq:schur} being absent because $t_z=0$.}
  \label{fig:boundary}
\end{figure}

\subsection{The principle}\label{sec:variational}

Corollary~\ref{cor:interiorcex} speaks about a given chart point.  The
closed statement and the statement about designs are the same statement,
so the hypotheses of Theorem~\ref{thm:boundary} are ordinary smaller
instances of the problem.

\begin{proposition}[Lean]\label{prop:density}
Write $\Gtz^{\mathrm{ch}}_{\circ}(m,k)$ for the assertion that every
chart point of size $m$ and rank $k$ with strictly positive weights
admits a dominating $k$-subset.  Then
$\Gtz^{\mathrm{ch}}_{\circ}(m,k)$ implies $\Gtz^{\mathrm{ch}}(m,k)$, at
the same size and rank.
\end{proposition}

\begin{proof}
Let $(P,t)$ be a chart point at which no $k$-subset dominates.  For each
of the finitely many $k$-subsets $C$ choose $x^{C}$ supported on $C$ with
$(x^{C})\tp Wx^{C}<0$.  The form is quadratic, so rescaling $x^{C}$ by
$u>0$ multiplies its value by $u^{2}$, and $x^{C}$ may be normalized to
$(x^{C})\tp Wx^{C}=-1$.  The common value lets one mixing parameter
serve every selection.  Put
$\mu_C:=\sum_c(1/m-t_c)(x^{C}_c)^{2}$ and
$\theta:=\bigl(1+\sum_C\lvert\mu_C\rvert\bigr)^{-1}\in(0,1]$.  By
\eqref{eq:relax},
\[
  (x^{C})\tp W^{\theta}x^{C}\;=\;-1-\theta\mu_C\;<\;0
  \qquad\text{for every }C,
\]
since
$\theta\lvert\mu_C\rvert\le\theta\sum_{C'}\lvert\mu_{C'}\rvert<1$, each
term being at most the whole sum.  So no $k$-subset dominates at
$(P,t^{\theta})$, whose weights are strictly positive by
Lemma~\ref{lem:relax}, contradicting
$\Gtz^{\mathrm{ch}}_{\circ}(m,k)$.
\end{proof}

\begin{corollary}\label{cor:equivalent}
The three statements $\Gtz(m,k)$, $\Gtz^{\mathrm{ch}}_{\circ}(m,k)$ and
$\Gtz^{\mathrm{ch}}(m,k)$ are equivalent, for every $m$ and $k$.
\end{corollary}

\begin{proof}
By Proposition~\ref{prop:chartdesign} the designs of size $m$ and rank
$k$ correspond to the chart points with positive weights, and domination
corresponds to domination; so $\Gtz(m,k)$ and
$\Gtz^{\mathrm{ch}}_{\circ}(m,k)$ are the same assertion.
Proposition~\ref{prop:density} gives
$\Gtz^{\mathrm{ch}}_{\circ}(m,k)\Rightarrow\Gtz^{\mathrm{ch}}(m,k)$, and
the converse is trivial, a chart point with positive weights being a
chart point.
\end{proof}

\begin{theorem}\label{thm:variational}
Let $m>k\ge1$ and assume $\Gtz(m-1,k)$ and $\Gtz(m-1,k-1)$.  If
$\Gtz(m,k)$ fails, then $G$ attains its minimum over the compact space of
chart points of size $m$ and rank $k$ at a point $(P_{\min},t_{\min})$
with
\[
  G(P_{\min},t_{\min})<0
  \qquad\text{and}\qquad
  (t_{\min})_c>0\ \text{ for every }c .
\]
In particular $(P_{\min},t_{\min})$ is the chart point of a design
$\dsn_{\min}$ of size $m$ and rank $k$; no $k$-subset dominates in
$\dsn_{\min}$; and $\dsn_{\min}$ minimizes $G$ over all designs of size $m$ and
rank $k$.
\end{theorem}

\begin{proof}
By Corollary~\ref{cor:equivalent} the hypotheses are
$\Gtz^{\mathrm{ch}}(m-1,k)$ and $\Gtz^{\mathrm{ch}}(m-1,k-1)$, so
Theorem~\ref{thm:boundary} and Corollary~\ref{cor:interiorcex} are
available at size $m$ and rank $k$.

A counterexample to $\Gtz(m,k)$ is a design in which no $k$-subset
dominates; by Proposition~\ref{prop:chartdesign} its chart point admits
no dominating $k$-subset, so $G<0$ there by \eqref{eq:threshold}, and the
minimum of $G$ is negative.  By Lemmas~\ref{lem:compact}
and~\ref{lem:continuity} that minimum is attained, at
$(P_{\min},t_{\min})$ say.  Then $G(P_{\min},t_{\min})<0$, so no
$k$-subset dominates at $(P_{\min},t_{\min})$, so $(t_{\min})_c>0$ for
every $c$ by Corollary~\ref{cor:interiorcex}, and $(P_{\min},t_{\min})$
is the chart point of a design $\dsn_{\min}$ by
Proposition~\ref{prop:chartdesign}.
Finally $\min G$ equals the infimum of $G$ over designs, by
Lemma~\ref{lem:relax}, so $\dsn_{\min}$ attains that infimum.
\end{proof}

At $(6,3)$ the hypotheses are $\Gtz(5,3)$ and $\Gtz(5,2)$, both
instances of Corollary~\ref{cor:corank}.  At $(7,3)$ they are
$\Gtz(6,3)$ and $\Gtz(6,2)$; the second is Theorem~\ref{thm:ranktwo},
and the first is the cell $(6,3)$ itself.
The two open cells of Corollary~\ref{cor:twocells} are therefore treated
in order: $(6,3)$ first, and its resolution supplies the hypothesis
needed at $(7,3)$.

By Corollary~\ref{cor:equivalent} the statement about chart points with
positive weights is the original problem.  The mixing of
Lemma~\ref{lem:relax} does move a counterexample into the interior, but
it destroys whatever distinguished the original point, so it cannot be
applied to a minimizer.  Minimizing over the design space directly fails
as well: that space is not compact, and a minimizing sequence of designs
need not converge to a design (Example~\ref{ex:escape}).  The minimum
exists only on the compactification, and
Corollary~\ref{cor:interiorcex} puts it back inside.

\begin{wip}
Transport stationarity across the congruence $D[C]^{1/2}$ of
Remark~\ref{rem:objectives}, or restate the system of
\S\ref{sec:interior} for $G$; either removes the hypothesis of
Definition~\ref{def:extremalcex} from Theorem~\ref{thm:rankthree}.
\end{wip}

\subsection{Sharpness of the boundary step}\label{sec:sharp}

Theorem~\ref{thm:schurstep} requires $t_z=0$ exactly.  Call the
\emph{compression lift} the composite of Lemma~\ref{lem:deflate},
Lemma~\ref{lem:delete} and Theorem~\ref{thm:schurstep}: deflate at $z$,
delete $z$, renormalize, lift a dominating subset back.  It is meaningful
at any $t_z<1$.  A minimizing-sequence argument would need the
compression lift and not the boundary statement, since along such a
sequence the weight at the escaping label is positive and merely tends to
zero.  The lift fails at every positive weight, and no threshold in $t_z$
repairs it.

\begin{lemma}\label{lem:pairdet}
Let $(P,t)$ be a chart point, let $z$ satisfy $P_{zz}>0$ and $t_z<1$, and
let $c\ne z$.  Let $P'$ be the deflation \eqref{eq:deflate} and let
$\varepsilon:=P'_{cc}-t_c/(1-t_z)$ be the gap of the compressed point at $c$.
Then
\begin{equation}\label{eq:family}
  \det W[\{c,z\}]
  \;=\;\varepsilon\,P_{zz}\;-\;t_z\Bigl(P_{cc}-t_c-\frac{t_cP_{zz}}{1-t_z}\Bigr) .
\end{equation}
\end{lemma}

\begin{proof}
By \eqref{eq:deflate}, $P'_{cc}=P_{cc}-P_{cz}^{2}/P_{zz}$, so
\[
  P_{cz}^{2}\;=\;P_{zz}\bigl(P_{cc}-P'_{cc}\bigr)
  \;=\;P_{zz}\Bigl(P_{cc}-\frac{t_c}{1-t_z}-\varepsilon\Bigr).
\]
Since $W[\{c,z\}]=
\left(\begin{smallmatrix}P_{cc}-t_c&P_{cz}\\P_{cz}&P_{zz}-t_z\end{smallmatrix}\right)$,
\begin{align*}
  \det W[\{c,z\}]
  &=(P_{cc}-t_c)(P_{zz}-t_z)
     -P_{zz}\Bigl(P_{cc}-\frac{t_c}{1-t_z}-\varepsilon\Bigr)\\
  &=P_{cc}P_{zz}-t_zP_{cc}-t_cP_{zz}+t_zt_c
     -P_{zz}P_{cc}+\frac{t_cP_{zz}}{1-t_z}+\varepsilon P_{zz}\\
  &=\varepsilon\,P_{zz}-t_z\Bigl(P_{cc}-t_c-\frac{t_cP_{zz}}{1-t_z}\Bigr),
\end{align*}
using $t_cP_{zz}/(1-t_z)-t_cP_{zz}=t_z\,t_cP_{zz}/(1-t_z)$ in the last
step.
\end{proof}

Identity \eqref{eq:family} separates the two contributions.
The compressed slack $\varepsilon$ enters undamped, with coefficient $P_{zz}>0$;
the weight at the pivot enters linearly, with a coefficient tending to
$P_{cc}-t_c(1+P_{zz})$ as $t_z\to0$.  Hold the chart and the survivor
fixed.  If $\varepsilon>0$ then the determinant tends to
$\varepsilon\,P_{zz}>0$, and the pair ceases to obstruct for all small
$t_z$.  If $\varepsilon=0$ then the determinant is
$-t_z(P_{cc}-t_c-t_cP_{zz}/(1-t_z))$, negative for every $t_z>0$ at which
the bracket is positive.

The lift therefore fails wherever the compressed slack vanishes against
a positive bracket, however small $t_z$ is.

\begin{example}\label{ex:liftfailure}
On the index set $\{0,1,2,3\}$ let
\begin{equation}\label{eq:liftchart}
  P\;=\;\frac13\begin{pmatrix}
    2&1&0&1\\ 1&1&1&0\\ 0&1&2&-1\\ 1&0&-1&1\end{pmatrix},
\end{equation}
and let the weights run over the one-parameter family
\begin{equation}\label{eq:liftweights}
  t_1=\tau,\quad t_0=t_2=t_3=\frac{1-\tau}{3}
  \qquad(0\le\tau\le1),
\end{equation}
whose member at $\tau=\tfrac14$ is the uniform vector.  Let $z=1$ be the
pivot and $c=0$ the survivor.  Write $M:=3P$.  The
ten inner products of the rows $m_0,\dots,m_3$ of $M$ are
\[
  \begin{array}{c|cccc}
   \langle m_i,m_j\rangle & j=0&1&2&3\\ \hline
   i=0 & 6&3&0&3\\
   1 & 3&3&3&0\\
   2 & 0&3&6&-3\\
   3 & 3&0&-3&3
  \end{array}
\]
that is, $M^{2}=3M$; hence $P^{2}=P$, and $P$ is symmetric with
$\tr P=(2+1+2+1)/3=2$.  So $(P,t)$ is a chart point of size $4$ and rank
$2$ for every $\tau$.

At the pivot, $P_{11}=\tfrac13$ and
$Pe_1=(\tfrac13,\tfrac13,\tfrac13,0)\tp$, so
$\lVert Pe_1\rVert^{2}=\tfrac13=P_{11}$, in accordance with
\eqref{eq:pivot}.  Deflation \eqref{eq:deflate} subtracts $\tfrac13$ from
every entry indexed in $\{0,1,2\}$:
\[
  P'\;=\;\frac13\begin{pmatrix}
    1&0&-1&1\\ 0&0&0&0\\ -1&0&1&-1\\ 1&0&-1&1\end{pmatrix},
  \qquad \tr P'=1 .
\]
Deleting the index $1$ leaves the three survivors with weights
$\bigl((1-\tau)/3\bigr)/(1-\tau)=\tfrac13$ each, independently of $\tau$.
The compressed gap at the survivor $0$ is therefore
\[
  \varepsilon\;=\;P'_{00}-\tfrac13\;=\;\tfrac13-\tfrac13\;=\;0
  \qquad\text{at every }\tau ,
\]
so $\{0\}$ dominates the compressed point --- with slack exactly zero, at
every parameter value.  Upstairs, by \eqref{eq:family} with $\varepsilon=0$,
$P_{zz}=\tfrac13$, $P_{cc}=\tfrac23$, $t_c=(1-\tau)/3$ and $t_z=\tau$,
\[
  \det W[\{0,1\}]
  \;=\;-\tau\Bigl(\frac23-\frac{1-\tau}{3}
      -\frac{1-\tau}{3}\cdot\frac{1/3}{1-\tau}\Bigr)
  \;=\;-\tau\cdot\frac{6-3(1-\tau)-1}{9}
  \;=\;-\frac{\tau(2+3\tau)}{9},
\]
negative for every $\tau>0$.  Explicitly
$W[\{0,1\}]=\left(\begin{smallmatrix}(1+\tau)/3&1/3\\
1/3&(1-3\tau)/3\end{smallmatrix}\right)$, and the probe
$x=(-1,1+\tau,0,0)$ gives
\begin{align*}
  x\tp Wx
  &\;=\;\frac{1+\tau}{3}-\frac{2(1+\tau)}{3}
      +\frac{(1-3\tau)(1+\tau)^{2}}{3}\\
  &\;=\;\frac{1+\tau}{3}\bigl(-1+(1-3\tau)(1+\tau)\bigr)
  \;=\;-\frac{(1+\tau)\tau(2+3\tau)}{3}.
\end{align*}
At $\tau=\tfrac14$ the block, its determinant and the probe $(1,-4)$ are
\[
  W[\{0,1\}]=\begin{pmatrix}5/12&1/3\\ 1/3&1/12\end{pmatrix},
  \qquad \det=-\tfrac{11}{144},
  \qquad (1,-4)\,W[\{0,1\}]\,(1,-4)\tp=-\tfrac{11}{12}.
\]
The leverages $P_{cc}/t_c$ at $\tau=\tfrac14$ are
$\tfrac83,\tfrac43,\tfrac83,\tfrac43$, all exceeding $1$, so the point is not
disposed of by deletion of a light atom.
\end{example}

\begin{wip}
Exhibit a one-parameter deformation of \eqref{eq:liftchart} with compressed
slack $\varepsilon>0$, and display the resulting determinant
$(3\varepsilon-2\tau-3\tau^{2})/9$ of \eqref{eq:family} changing sign at
$\tau\asymp \varepsilon$.
\end{wip}

\begin{theorem}[Lean]\label{thm:sharp}
There is no $\tau_0>0$ such that the compression lift holds at every
chart point with $t_z<\tau_0$.  Explicitly: for every $\tau_0>0$ there
are a chart point of size $4$ and rank $2$, a pivot $z$ with $P_{zz}>0$
and $0<t_z<\tau_0$, and a subset dominating the compressed point of size
$3$ and rank $1$, whose image together with $z$ does not dominate.
\end{theorem}

\begin{proof}
Take the family of Example~\ref{ex:liftfailure} at
$\tau=\min\{\tau_0,1\}/2$.  Then $0<\tau<\tau_0$ and $\tau<1$; the
compressed point is dominated by $\{0\}$, since its gap there is $0$; and
$\{0,1\}$ does not dominate upstairs, the probe $(-1,1+\tau,0,0)$ giving
$-(1+\tau)\tau(2+3\tau)/3<0$.
\end{proof}

The obstruction sits on the equality locus of \S\ref{sec:ties}.
Shrinking $t_z$ does shrink the deficit, but along this family it shrinks
nothing else: the compressed slack is frozen at zero, so the deficit is
never overtaken.  By \eqref{eq:family} a positive compressed slack would
overtake it, and no positive lower bound on $\varepsilon$ is to be had.
The compressed point is a chart point one size and one rank lower, hence
by Proposition~\ref{prop:chartdesign} the chart point of a design
wherever its weights are positive; and by
Theorem~\ref{thm:tieseverywhere} designs whose dominating subsets all
have $\lmin(\SC)=1$, that is with chart slack zero, exist at every cell
with $m\ge k+1$ and at every weight vector.

The compression sends the whole family to one point, the surviving
weights being $\tfrac13$ at every $\tau$.  That point charts the design
$g=(1,-1,1)$ at $(3,1)$, of leverage $1$ at each index, so each of its
singletons dominates with $\lmin(\SC)=1$.
Figure~\ref{fig:boundarytrace} runs the four steps and the fork on the
family.

\begin{figure}[ht]
  \centering
  % The boundary dichotomy executed on the chart point of ex:liftfailure at (4,2),
% with the intermediate matrix printed at every step, and the fork at the lift.
%
% PROJECTION.  There is none: the figure is a flow of text and matrices, and no
% space of dimension above zero is drawn.  The layout is a grid in picture
% centimetres, and every ordinate below is a node CENTRE, set so that
% consecutive boxes clear one another by about 9pt, which is twice the 1.7mm
% arrowhead that joins them.  The datum spans the top from the ordinate 9.15
% down.  The four steps of the compression are stacked in one column, centred at
% the abscissa 3.45 on the ordinates 6.52, 4.57, 2.22 and -0.01 under a title on
% 7.80, and joined downwards by arrows; they are the steps that do not see the
% parameter.  The stem at the abscissa 7.35 carries the compressed verdict into
% both branches, which are centred at the abscissa 11.35 on the ordinates 4.95
% and 1.93, with the identity they share above them on 7.80, a rule on 0.12, and
% the standing pair anchored north on -0.06.  Every number in the body is an
% entry of one of the matrices derived below, and no drawn coordinate carries
% data.
%
% THE CHART POINT, exactly, as in eq:liftchart and eq:liftweights.
%   P = (1/3) [[2,1,0,1],[1,1,1,0],[0,1,2,-1],[1,0,-1,1]],
%   t = ( (1-tau)/3, tau, (1-tau)/3, (1-tau)/3 ),   0 <= tau <= 1 .
% With M := 3P the ten inner products of the rows of M are
%   <m_0,m_0> = 6  <m_0,m_1> = 3  <m_0,m_2> = 0  <m_0,m_3> =  3
%   <m_1,m_1> = 3  <m_1,m_2> = 3  <m_1,m_3> = 0  <m_2,m_2> =  6
%   <m_2,m_3> = -3 <m_3,m_3> = 3
% that is M^2 = 3M, whence P^2 = P; P is symmetric and tr P = (2+1+2+1)/3 = 2.
% The weights are nonnegative and sum to 3*(1-tau)/3 + tau = 1.  So (P,t) is a
% chart point of size four and rank two at every tau in [0,1].
%
% THE RANGE.  ex:liftfailure declares 0 <= tau <= 1 and eq:liftweights the open
% interval, which is where the lift fails and where the chart point is the chart
% point of a design.  The figure takes the closed range: the left branch of the
% fork is the member at tau = 0, the one thm:schurstep applies to, and that
% member is a chart point and not a design.  Lean's Gtz.liftFailurePointAt
% carries the same closed range, 0 <= pivotWeight <= 1, and its frozen-slack
% theorem carries in addition the tau < 1 that STEP 3 needs.
%
% STEP 1, the pivot, at z = 1.  Column 1 of P is (1,1,1,0)/3, of squared length
% 3*(1/9) = 1/3, which is P_11: this is eq:pivot, and P_11 > 0 puts the point in
% the second branch of thm:boundary.
%
% STEP 2, deflation.  (P e_1)(P e_1)^T / P_11 = 3 * (1/9) J on the index block
% {0,1,2}, that is 1/3 at each of the nine entries indexed there and 0 elsewhere,
% so eq:deflate subtracts 1/3 from every entry indexed in {0,1,2}:
%   P' = (1/3) [[1,0,-1,1],[0,0,0,0],[-1,0,1,-1],[1,0,-1,1]] .
% With u = (1,0,-1,1) one has u u^T equal to that integer matrix and ||u||^2 = 3,
% so P' = u u^T/||u||^2 is the orthogonal projection on the line R u: symmetric,
% idempotent, of trace 1 = k - 1, with P' e_1 = 0.
%
% STEP 3, deletion.  Striking the row and column 1 leaves, on the survivors
% {0,2,3} in that order,
%   hat P = (1/3) [[1,-1,1],[-1,1,-1],[1,-1,1]] ,
% and lem:delete, whose hypothesis is t_z < 1, renormalizes each surviving weight
% to ((1-tau)/3)/(1-tau) = 1/3, the same number at every tau < 1: the compression
% sends the whole one-parameter family to one point.  This is the one step of the
% four that does not run at tau = 1.  That point has positive weights,
% so by prop:chartdesign it is the chart point of a design at (3,1); the rows of
% V are the entries of u/sqrt3 = (1,-1,1)/sqrt3, and g_c = v_c/sqrt(t_c) gives
%   g = (1,-1,1),   ell_c = 1,   sum_c t_c g_c^2 = 3*(1/3)*1 = 1 .
% Every singleton has S_C = 1, so every singleton dominates with equality: the
% compressed point lies on the equality locus of section sec:ties.
%
% STEP 4, the compressed gap.  hat W = hat P - diag(1/3,1/3,1/3) =
% (1/3)[[0,-1,1],[-1,0,-1],[1,-1,0]], whose diagonal is zero.  So the compressed
% slack of lem:pairdet is eps = hat W_00 = 0, at every tau.
%
% THE IDENTITY.  For x supported on {0,1}, write x = y + x_1 e_1 with y = x_0 e_0.
% Then <P e_1, y> = x_0/3, P'_00 - t_0 = 1/3 - (1-tau)/3 = tau/3, and eq:schur is
%   x^T W x = (tau/3) x_0^2 + (1/3)(x_0 + x_1)^2 - tau x_1^2 ,
% which expands to ((1+tau)/3)x_0^2 + (2/3)x_0 x_1 + (1/3 - tau)x_1^2, the form of
% W[{0,1}].  The three terms are the deflated gap, the square and the deficit.
%
% THE FORK.  At tau = 0 the deficit is absent and the first term vanishes with
% it, leaving (1/3)(x_0+x_1)^2 >= 0; there W[{0,1}] = (1/3)[[1,1],[1,1]], of trace
% 2/3 and determinant 0, hence of spectrum {0, 2/3}.  At tau > 0,
%   W[{0,1}] = [[(1+tau)/3, 1/3],[1/3, 1/3 - tau]],
%   det = ((1+tau)(1-3tau) - 1)/9 = -tau(2+3tau)/9 ,
% which is eq:family at eps = 0, P_zz = 1/3, P_cc = 2/3, t_c = (1-tau)/3; and the
% probe (-1, 1+tau, 0, 0) gives ((1+tau)/3)(1 - 2 + (1+tau)(1-3tau)) =
% -tau(1+tau)(2+3tau)/3.  At tau = 1/4 the block is [[5/12,1/3],[1/3,1/12]] of
% determinant 5/144 - 16/144 = -11/144, and (1,-4,0,0) gives
% 5/12 - 8/3 + 16/12 = -11/12.
%
% THE STANDING PAIR.  P_02 = 0, P_00 = P_22 = 2/3 and t_0 = t_2 = (1-tau)/3, so
% W[{0,2}] = ((1+tau)/3) I_2, positive semidefinite at every tau in [0,1].  The
% pair {0,2} therefore dominates throughout the family.
%
% DISCLOSED.  The Lean development carries the chart (symmetric, idempotent, of
% trace two), the one-parameter family, the frozen compressed slack and the
% failure of the lift at every positive weight, together with the three general
% steps; it does not state the displayed P', hat P and hat W, which are computed
% inside its proofs, and it does not state the standing pair {0,2}.  The passage
% from the compressed chart point to the design g = (1,-1,1) uses the half of
% prop:chartdesign that section sec:formalization records as not mechanized.
\begin{tikzpicture}[
  bx/.style={draw, rounded corners=1pt, line width=0.5pt, inner sep=4pt,
             align=center, font=\scriptsize},
  ar/.style={-{Latex[length=1.7mm]}, line width=0.5pt},
  ttl/.style={font=\scriptsize\itshape, align=center},
  txt/.style={font=\scriptsize, align=center, inner sep=1.5pt},
  hair/.style={line width=0.4pt, black!45}]

% ====================== the datum, spanning the top =======================
\node[anchor=north west, align=left, font=\scriptsize, inner sep=0pt]
  at (0,9.15) {%
  $P=\tfrac13\begin{psmallmatrix}2&1&0&1\\1&1&1&0\\0&1&2&-1\\
                                 1&0&-1&1\end{psmallmatrix}$, \
  $t^{(\tau)}=\bigl(\tfrac{1-\tau}{3},\ \tau,\ \tfrac{1-\tau}{3},
                    \ \tfrac{1-\tau}{3}\bigr)$, \
  $M:=3P$ has $M^{2}=3M$, so $P^{2}=P$; and $P\tp=P$, $\tr P=2$\\[3pt]
  a chart point of size four and rank two at every $\tau\in[0,1]$
  (Definition~\ref{def:chartpoint}), with escaping index $z=1$};

% ================ the compression, four steps, one column =================
\node[ttl, anchor=north] at (3.45,7.80) {the four steps below are free of
  $\tau$};

\node[bx] (s1) at (3.45,6.52) {%
  \textit{the pivot} (Lemma~\ref{lem:pivot})\\[3pt]
  $Pe_1=\tfrac13(1,1,1,0)\tp$, \
  $\lVert Pe_1\rVert^{2}=\tfrac13=P_{11}>0$};

\node[bx] (s2) at (3.45,4.57) {%
  \textit{deflation} (Lemma~\ref{lem:deflate})\\[3pt]
  $P'=\tfrac13\begin{psmallmatrix}1&0&-1&1\\0&0&0&0\\-1&0&1&-1\\
                                  1&0&-1&1\end{psmallmatrix}
      =\tfrac13uu\tp$, \ $u=(1,0,-1,1)$\\[3pt]
  $\tr P'=1$, \ $P'e_1=0$};

\node[bx] (s3) at (3.45,2.22) {%
  \textit{deletion} (Lemma~\ref{lem:delete}, at $\tau<1$)\\[3pt]
  $\widehat P=\tfrac13\begin{psmallmatrix}1&-1&1\\-1&1&-1\\
                                          1&-1&1\end{psmallmatrix}$
  on $\{0,2,3\}$, \
  $\widehat t_c=\dfrac{(1-\tau)/3}{1-\tau}=\tfrac13$\\[3pt]
  its design at $(3,1)$: \ $g=(1,-1,1)$, \ $\ell_c=1$};

\node[bx] (s4) at (3.45,-0.01) {%
  \textit{the compressed point}\\[3pt]
  $\widehat W=\tfrac13\begin{psmallmatrix}0&-1&1\\-1&0&-1\\
                                          1&-1&0\end{psmallmatrix}$, \
  $\varepsilon=\widehat W_{00}=0$\\[3pt]
  $\{0\}$ dominates, with slack $0$};

\draw[ar] (s1) -- (s2);
\draw[ar] (s2) -- (s3);
\draw[ar] (s3) -- (s4);

% ================= the identity that both branches read ===================
\node[txt, anchor=north] at (11.35,7.80) {%
  for $x$ supported on $\{0,1\}$, identity \eqref{eq:schur} reads\\[3pt]
  $x\tp Wx=\tfrac{\tau}{3}x_0^{2}
     +\tfrac13\bigl(x_0+x_1\bigr)^{2}-\tau x_1^{2}$\\[3pt]
  the deflated gap, a square, and the deficit};

% ================================ the fork ================================
\node[bx] (f0) at (11.35,4.95) {%
  \textit{the lift, at} $\tau=0$\\[3pt]
  only the square survives: \
  $x\tp Wx=\tfrac13(x_0+x_1)^{2}\ge0$\\[3pt]
  $W[\{0,1\}]=\tfrac13\begin{psmallmatrix}1&1\\1&1\end{psmallmatrix}$, \
  $\operatorname{spec}=\{0,\tfrac23\}$\\[3pt]
  Theorem~\ref{thm:schurstep}: $\{0,1\}$ dominates, $\lmin=0$};

\node[bx] (f1) at (11.35,1.93) {%
  \textit{the lift fails, at} $\tau>0$\\[3pt]
  $W[\{0,1\}]=\begin{psmallmatrix}(1+\tau)/3&1/3\\
                                  1/3&1/3-\tau\end{psmallmatrix}$, \
  and \eqref{eq:family}\\[3pt]
  at $\varepsilon=0$ gives $\det=-\tfrac{\tau(2+3\tau)}{9}$\\[3pt]
  $x=(-1,1+\tau,0,0)$ gives $-\tfrac13\tau(1+\tau)(2+3\tau)$\\[3pt]
  at $\tau=\tfrac14$: \ $\det=-\tfrac{11}{144}$ and
  $(1,-4,0,0)$ gives $-\tfrac{11}{12}$};

\draw[line width=0.5pt] (s4.east) -| (7.35,4.95);
\draw[ar] (7.35,4.95) -- (f0.west);
\draw[ar] (7.35,1.93) -- (f1.west);

% ========================== what stands throughout ========================
\draw[hair] (7.80,0.12) -- (15.00,0.12);
\node[txt, anchor=north] at (11.35,-0.06) {%
  $W[\{0,2\}]=\tfrac{1+\tau}{3}I_2\psd0$ at every $\tau$:\\[2pt]
  the pair $\{0,2\}$ dominates throughout};

\end{tikzpicture}

  \caption{The dichotomy of Theorem~\ref{thm:boundary} executed on the
  chart point of Example~\ref{ex:liftfailure}, at $(4,2)$ with pivot $z=1$.
  Left: the four steps of the compression lift, each with the matrix it
  produces; none of them sees $\tau$, and the deletion returns the same
  compressed point at every $\tau$, the chart point of the design
  $g=(1,-1,1)$ at $(3,1)$ on which every singleton dominates with
  $\lmin(\SC)=1$.  Right: identity \eqref{eq:schur} on the pair
  $\{0,1\}$, and the two branches it splits into.  At $\tau=0$ the
  deflated gap and the deficit vanish together and
  Theorem~\ref{thm:schurstep} lifts $\{0\}$ to $\{0,1\}$ with least
  eigenvalue $0$, so the lift leaves no slack.  At $\tau>0$ the deficit
  survives, and \eqref{eq:family} at $\varepsilon=0$ turns it into a
  negative determinant.  The pair $\{0,2\}$ dominates throughout the
  family, so what fails at positive weight is the lift rather than the
  chart point.}
  \label{fig:boundarytrace}
\end{figure}

This is Corollary~\ref{cor:nomargin} again, in local form.  It is why the
boundary is discharged at $t_z=0$ by the identity \eqref{eq:schur}
instead of at small $t_z$ by an estimate, and why
Theorem~\ref{thm:variational} speaks of the limit point rather than of a
minimizing sequence.  Figure~\ref{fig:sharp} shows the two contributions.

\begin{figure}[ht]
  \centering
  % Identity (eq:family) evaluated on the chart of section app:liftfail:
%   P_zz = 1/3, P_cc = 2/3, t_c = (1-t)/3, t_z = t,
%   det W[{c,z}] = d/3 - t(2+3t)/9 = (3d - 2t - 3t^2)/9 .
% Level v is the parabola  d = 3v + 2t/3 + t^2 ; the zero level is
%   d = 2t/3 + t^2 , tangent at the origin to the line d = 2t/3.
\begin{tikzpicture}
\begin{axis}[
  name=sharpaxis,
  width=8.4cm, height=6.6cm,
  scale only axis,
  xmin=0, xmax=0.3, ymin=0, ymax=0.3,
  xtick={0,0.05,0.1,0.15,0.2,0.25,0.3},
  ytick={0,0.05,0.1,0.15,0.2,0.25,0.3},
  xticklabel style={font=\scriptsize, /pgf/number format/fixed,
                    /pgf/number format/precision=2},
  yticklabel style={font=\scriptsize, /pgf/number format/fixed,
                    /pgf/number format/precision=2},
  xlabel={$t_z$}, ylabel={$\varepsilon$},
  xlabel style={font=\small}, ylabel style={font=\small, rotate=-90},
  samples=200, domain=0:0.3,
  clip=true]

% ---- the region where the lift survives -----------------------------
\fill[black!12] (axis cs:0,0) rectangle (axis cs:0.3,0.3);
\addplot[draw=none, fill=white] {2*x/3 + x^2} \closedcycle;

% ---- level curves  (3d - 2t - 3t^2)/9 = v ---------------------------
\addplot[black!55, line width=0.35pt, densely dotted] {2*x/3 + x^2 - 0.18};
\addplot[black!55, line width=0.35pt, densely dotted] {2*x/3 + x^2 - 0.12};
\addplot[black!55, line width=0.35pt, densely dotted] {2*x/3 + x^2 - 0.06};
\addplot[black!55, line width=0.35pt, dashed] {2*x/3 + x^2 + 0.06};
\addplot[black!55, line width=0.35pt, dashed] {2*x/3 + x^2 + 0.12};
\addplot[black!55, line width=0.35pt, dashed] {2*x/3 + x^2 + 0.18};

% ---- the tangent line d = 2t/3 at the origin ------------------------
\addplot[black!70, line width=0.35pt, dash pattern=on 1pt off 1.2pt]
  {2*x/3};

% ---- the zero level --------------------------------------------------
\addplot[black, line width=1.1pt] {2*x/3 + x^2};

% ---- the extremal axis d = 0 ----------------------------------------
\draw[line width=1.4pt] (axis cs:0,0) -- (axis cs:0.3,0);

% ---- labels inside the frame ----------------------------------------
\node[font=\scriptsize, anchor=north west] at (axis cs:0.006,0.294)
  {lift survives};
\node[font=\scriptsize, anchor=south east] at (axis cs:0.294,0.014)
  {lift fails};
\draw[line width=0.3pt] (axis cs:0.150,0.1225) -- (axis cs:0.183,0.150);
\node[font=\scriptsize, anchor=west, fill=white, inner sep=1pt]
  at (axis cs:0.185,0.152) {zero level};
\node[font=\scriptsize, rotate=25, fill=white, inner sep=1pt]
  at (axis cs:0.262,0.176) {$\varepsilon=\tfrac23 t_z$};
\node[font=\scriptsize, anchor=south west, fill=white, inner sep=1pt]
  at (axis cs:0.040,0.003) {$\varepsilon=0$: the value is $-t_z(2+3t_z)/9<0$};
\end{axis}

% ---- the reading of the picture, below the frame ---------------------
\node[font=\scriptsize, anchor=north, align=center]
  at ($(sharpaxis.south)+(0,-0.95cm)$)
  {On the chart of Example~\ref{ex:liftfailure} identity
   \eqref{eq:family} reads
   $\det W[\{c,z\}]=\bigl(3\varepsilon-2t_z-3t_z^{2}\bigr)/9$,\\
   with zero level $\varepsilon=\tfrac23t_z+t_z^{2}$.
   Dotted: the levels $-\tfrac1{50},-\tfrac1{75},-\tfrac1{150}$.
   Dashed: their positives.};
\end{tikzpicture}

  \caption{Sharpness of the compression lift.  The compressed slack $\varepsilon$
  enters \eqref{eq:family} undamped, the pivot weight $t_z$ linearly.  The
  family of Example~\ref{ex:liftfailure} is frozen on the axis $\varepsilon=0$.  The
  zero level leaves the origin with slope $\tfrac23$, so that axis lies
  strictly inside the failure region for every $t_z>0$, and no threshold
  in $t_z$ repairs the lift.}
  \label{fig:sharp}
\end{figure}
